\begin{helper}
Let \(0<\eta<1/16\).  For all sufficiently large \(n\), with
\[
L=\log n,\qquad
k_R=(1+\eta)\sqrt{nL},\qquad
K=\noderef{TwoBiteNaturalIndependenceScale}(1+\eta,n),\qquad
k=(K:\mathbb R),
\]
and \(m_R=n/L^2\), one has
\[
(2kL+L)\bigl(16Lm_R n^{8\eta}\bigr)\le k_R^4 .
\]
\end{helper}
\begin{proof}
Since \(K=\lceil k_R\rceil\), once \(k_R\ge1\) we have
\[
k\le k_R+1\le2k_R .
\]
Also, after increasing the threshold, \(L\ge1\).  Hence
\[
2kL+L\le4k_RL+L\le5k_RL .
\]
Substituting \(m_R=n/L^2\) gives
\[
(2kL+L)(16Lm_Rn^{8\eta})
\le 5k_RL\cdot16L\frac n{L^2}n^{8\eta}
=80k_Rn^{1+8\eta}.
\]
Using \(k_R=(1+\eta)\sqrt{nL}\), it remains enough to have
\[
80(1+\eta)n^{3/2+8\eta}\sqrt L
\le (1+\eta)^4n^2L^2 .
\]
Equivalently,
\[
80\le (1+\eta)^3 n^{1/2-8\eta}L^{3/2}.
\]
The exponent \(1/2-8\eta\) is positive because \(\eta<1/16\), and
\(L=\log n\) tends to infinity.  Thus the right-hand side tends to infinity,
so a final enlargement of the threshold makes the displayed inequality hold
for every \(n\) above the threshold.
\end{proof}
